\section{Problem formulation}
Let the structural parameters of client $i$ be $\vartheta_i$ and the measured
scheduling variable be $\rho_k$. During the oracle study,
\begin{equation}
  \vartheta_i(t)=\vartheta_i,\qquad \rho_k=v_{x,k}\in[10,30]\,\mathrm{m\,s^{-1}}.
\end{equation}
The eventual hierarchical model is
\begin{equation}
 \Theta_i(\rho)=\Theta_{\mathrm{shared}}(\rho)
 +\Delta\Theta_{c(i)}(\rho)+\epsilon_i(\rho),
 \qquad \Theta=[A\ B].
\end{equation}
The first finite-dimensional LPV basis to test is
\begin{equation}
 \Theta(\rho)=\Theta_0+\rho^{-1}\Theta_1+\rho^{-2}\Theta_2.
\end{equation}
Phase~1 must validate this approximation before it is used for learning.

The controlled output is yaw rate $r$, with sideslip $\beta$ constrained through
the cost. A two-degree-of-freedom scheduled LQI architecture is considered:
\begin{equation}
 u=-K_x(\rho)x-K_I(\rho)\chi+k_r(\rho)r^\star.
\end{equation}
The same $Q$ and $R$ weights are used for all methods so that improvements cannot
be attributed to method-specific tuning.
